\section{Risk Register}
\label{sec:risks}

\subsection{Method}

Each risk is scored on likelihood $L$ and impact $I$, both $1$--$5$, with
exposure $E = L \times I$. Risks with $E \geq 12$ require an owner, a named
mitigation, and a trigger condition that converts the mitigation into action.
Risks are reviewed at every phase boundary.

\subsection{Technical Risks}

\begin{table}[htbp]
\caption{Technical Risk Register}
\label{tab:techrisk}
\begin{center}
\renewcommand{\arraystretch}{1.15}
\begin{tabularx}{\columnwidth}{@{}l>{\raggedright\arraybackslash}Xccc@{}}
\toprule
\textbf{ID} & \textbf{Risk} & $L$ & $I$ & $E$ \\
\midrule
T-01 & Reference sensitometric data is unavailable or inconsistent for target stocks, preventing AC-1 validation. & 4 & 5 & 20 \\
T-02 & Model is numerically correct but perceptually unconvincing; discrimination rate stays high. & 3 & 5 & 15 \\
T-03 & Apple changes \texttt{CIRAWFilter} behavior so that scene-referred decode is no longer achievable. & 2 & 5 & 10 \\
T-04 & Performance targets unreachable on A14/A15, forcing a raised minimum device. & 2 & 4 & 8 \\
T-05 & Memory budget exceeded on 6\,GB devices during export, causing jetsam. & 3 & 4 & 12 \\
T-06 & Grain lattice produces visible aliasing at some format/resolution combinations. & 3 & 3 & 9 \\
T-07 & Profile version migration silently changes users' existing edits. & 2 & 5 & 10 \\
T-08 & Device calibration constant $g_{\mathrm{cal}}$ varies within a device model, breaking the exposure anchor. & 2 & 4 & 8 \\
T-09 & Halation pyramid produces visible seams at ROI tile boundaries under extreme radii. & 3 & 3 & 9 \\
T-10 & Thermal throttling makes sustained editing unpleasant on older devices. & 3 & 3 & 9 \\
\bottomrule
\end{tabularx}
\end{center}
\end{table}

\paragraph{T-01 (E = 20), owner: imaging scientist}
The highest-exposure technical risk. Published sensitometric data varies in
completeness; some stocks publish full $D$--$\logE$ families, others publish
only a summary. Digitizing a printed curve introduces its own error.

\emph{Mitigation.} Budget for physical measurement (Table~\ref{tab:cost}):
acquire film, expose calibrated step wedges under controlled conditions, process
at a professional laboratory with documented chemistry, and measure with a
transmission densitometer. This produces primary data. \emph{Trigger:} if by
month 4 fewer than six of the ten target stocks have data adequate for AC-1
validation, reduce the launch stock list to those that do and reallocate. A
smaller library of stocks that are demonstrably correct is a better product than
a larger library of guesses, and it is a better story.

\paragraph{T-02 (E = 15), owner: product lead}
It is entirely possible to build a model that matches every measurement and
still produces images that do not read as film, because the cues the eye uses may
not be the ones measured. This is the risk that the whole enterprise is
subtly aimed at the wrong target.

\emph{Mitigation.} The cue-elicitation protocol in
Section~\ref{sec:validation}: the panel names what tipped them off, and that list
drives development priority. \emph{Trigger:} if round-2 discrimination exceeds
80\% and the cue list is dominated by a phenomenon not in the model, escalate to
a model extension decision at the phase boundary rather than continuing to
refine existing stages.

\paragraph{T-05 (E = 12), owner: GPU engineer}
Table~\ref{tab:memory} shows an 878\,MB peak against a 900\,MB budget. That is
not enough margin.

\emph{Mitigation.} Export tiling is implemented in v1.0 rather than deferred,
and the tiling threshold is set conservatively --- tiling engages whenever the
projected live set exceeds 600\,MB, not 900\,MB. Additionally, the decoded
source is released before the encode staging buffer is allocated, which the
naive ordering does not do. \emph{Trigger:} any jetsam observed during Phase 2
device testing forces the tiling threshold down a further 100\,MB.

\subsection{Product and Market Risks}

\begin{table}[htbp]
\caption{Product Risk Register}
\label{tab:prodrisk}
\begin{center}
\renewcommand{\arraystretch}{1.15}
\begin{tabularx}{\columnwidth}{@{}l>{\raggedright\arraybackslash}Xccc@{}}
\toprule
\textbf{ID} & \textbf{Risk} & $L$ & $I$ & $E$ \\
\midrule
P-01 & The addressable audience for physical accuracy is smaller than modelled; conversion falls below 2\%. & 3 & 5 & 15 \\
P-02 & The laboratory vocabulary is a barrier rather than a differentiator for the P3 persona. & 3 & 4 & 12 \\
P-03 & A well-resourced competitor ships an adequate approximation first and defines the category. & 3 & 4 & 12 \\
P-04 & Apple ships comparable film simulation as a system feature. & 2 & 5 & 10 \\
P-05 & Trademark or brand objection to descriptive use of stock names. & 3 & 3 & 9 \\
P-06 & Subscription fatigue suppresses conversion regardless of quality. & 3 & 3 & 9 \\
\bottomrule
\end{tabularx}
\end{center}
\end{table}

\paragraph{P-01 (E = 15)}
Addressed in Section~\ref{sec:business}. The mitigation is early measurement
with a non-expert cohort during Phase 1, and the prepared response is
repositioning toward the professional tier with adjusted pricing, which the
architecture supports without rework.

\paragraph{P-02 (E = 12), owner: product designer}
The vocabulary is a core commitment, and abandoning it would hollow out the
product. But if a user cannot find the brightness control because it is called
print density, the commitment has failed in execution rather than in principle.

\emph{Mitigation.} The explanation affordance (long-press) is not optional
polish; it is the mechanism that makes the vocabulary teachable, and it is
scheduled in Phase 4 with dedicated content authoring. Additionally, search
within the develop workspace accepts colloquial terms and navigates to the
corresponding laboratory control, showing both names --- this teaches rather
than aliases. \emph{Trigger:} if beta task-completion testing shows more than
25\% of P3-persona users failing to locate a lightness control within 60
seconds, add a first-run guided tour before ship.

\paragraph{P-05 (E = 9), owner: product lead}
Film stock names are used descriptively to identify a modelled response, which
is a recognized nominative use, and no proprietary profile data is
redistributed (CR-04). Nonetheless the risk is nonzero.

\emph{Mitigation.} Legal review before launch (budgeted). A prepared fallback
naming scheme exists in which each profile carries a descriptive name
("Daylight Portrait 400") with the reference stock identified only in the
profile's technical notes. The fallback is designed now, so that adopting it
later is a data change and not an engineering change.

\subsection{Programme Risks}

\begin{table}[htbp]
\caption{Programme Risk Register}
\label{tab:progrisk}
\begin{center}
\renewcommand{\arraystretch}{1.15}
\begin{tabularx}{\columnwidth}{@{}l>{\raggedright\arraybackslash}Xccc@{}}
\toprule
\textbf{ID} & \textbf{Risk} & $L$ & $I$ & $E$ \\
\midrule
G-01 & Imaging scientist role cannot be filled, or is filled late. & 4 & 5 & 20 \\
G-02 & Phase 1 overruns, compressing everything downstream. & 4 & 4 & 16 \\
G-03 & Scope expansion into video before v1.0 ships. & 3 & 4 & 12 \\
G-04 & Key-person dependency on the GPU architecture. & 3 & 4 & 12 \\
\bottomrule
\end{tabularx}
\end{center}
\end{table}

\paragraph{G-01 (E = 20), owner: product lead}
The single highest-exposure risk in the programme. The required skill
combination is scarce.

\emph{Mitigation.} Begin recruiting before Phase 0. Structure the role to be
fillable by two adjacent profiles working together --- a color scientist without
strong implementation skills paired with an engineer without domain background
--- since that pair is far easier to hire than the unicorn. Engage the colorist
consultant early to provide domain review in the interim. \emph{Trigger:} if the
role is unfilled at month 2, activate the pair structure and extend Phase 1 by
six weeks in the plan of record rather than absorbing the slip silently.

\paragraph{G-02 (E = 16)}
Phase 1 is already given twice an optimistic estimate, which is the primary
mitigation. The secondary mitigation is the designated cut list in
Section~\ref{sec:roadmap} (optical, aging, inspector), which provides
approximately ten weeks of recoverable schedule in later phases without touching
any MUST requirement.

\paragraph{G-04 (E = 12)}
The render graph is subtle and, written by one person, becomes a bus factor of
one.

\emph{Mitigation.} The graph's design is documented in
Section~\ref{sec:rendergraph} at implementation detail, its properties are
enforced by the property-based tests of V-18 through V-20, and a second engineer
is assigned to Phase 2 explicitly for knowledge distribution rather than
throughput.

\subsection{Risk Summary}

\begin{table}[htbp]
\caption{Top Risks by Exposure}
\label{tab:toprisk}
\begin{center}
\renewcommand{\arraystretch}{1.15}
\begin{tabular}{@{}llc@{}}
\toprule
\textbf{ID} & \textbf{Risk} & $E$ \\
\midrule
T-01 & Reference data unavailable & 20 \\
G-01 & Imaging scientist hire & 20 \\
G-02 & Phase 1 overrun & 16 \\
T-02 & Perceptually unconvincing & 15 \\
P-01 & Audience smaller than modelled & 15 \\
\bottomrule
\end{tabular}
\end{center}
\end{table}

Four of the top five risks concern the imaging model or the people who build it.
That concentration is not a defect of the analysis; it is the accurate shape of
this programme, and it is the reason the schedule, the budget, and the staffing
plan all front-load the same area.
